\chapter[Sermon 21. The Lord Gives Wholesome Counsel to the Laodiceans, Admonishing Them to Repent.]{The Lord Gives Wholesome Counsel to the Laodiceans, Admonishing Them to Repent.}
\label{sermon:021}
\markright{Sermon 21. The Lord Gives Wholesome Counsel to the Laodiceans ...}

I advise you to buy from me gold refined by fire, so that you may be wealthy; and white clothing, so that you may be clothed, and the shame of your nakedness may not be revealed; and anoint your eyes with ointment, so that you may see. As many as I love, I rebuke and chasten. Be zealous, therefore, and repent.

\index{Repentance}Since God does not desire the death of a sinner, but rather that he should repent and live: Therefore, after severely rebuking the church of Laodicea, he gives her wholesome counsel, admonishing and exhorting her to repentance, and indicates what true repentance is.

The Lord uses the word of counseling, not of commanding, to attempt to confound the madness of those who, unless they are forcefully drawn, do not think themselves admonished, allured, or called by the Lord. And while they look for such a drawing, they neglect all of God’s counsel and fall from true salvation. God counsels his elect with things that are wholesome; the chosen obey good counsel. God touches their hearts inwardly, and outwardly by preaching of the word, and by sundry admonitions he pulls and drives man from evil to good. This counsel of God is not to be despised, and another violent vocation to be imagined. God’s word must be heard. “Today,” says the Prophet, “if you hear his voice, do not harden your hearts.” When the Lord counsels with his word, and the hearers harden their minds, they do so through their own fault and are made authors of their own destruction. But those who receive God’s counsel receive it not by the force of free will, but by the grace of God, which works in us to will and to perform.

Therefore, when the Lord advises beneficial things, the elect pray that they may receive them; and they receive them through grace, obeying the counsels of God.

And part of the wholesome counsel is this: “I know your works, says the Lord, that you have tried and have found Me faithful, as gold tried in the fire; that you may be rich, may be clothed, and may receive eye salve to anoint your eyes.” He sets these things as a medicine against the diseases which He revealed before, calling the church of the Laodiceans poor, naked, and blind. Now therefore He teaches them how they may be rich, may be clothed, and may receive their sight again, if they truly obtain gold tried in the fire, or are purified.

And gold tried in the fire is gold most purified and clean, having no base metals or impurities, but pure and clean gold. This shadows the word of God, of which the Prophet sang: “The word of the Lord is a pure word, silver tried in the fire, seven times purged in a vessel of earth.” Certainly, the word of God is light, coming from the eternal and most pure light, having no part of human defilement or affections, savoring of no errors, teaching nothing that is corrupt. However, of itself it will profit a man nothing unless it is received with a true and sincere faith. Therefore, I do not separate faith from the word, and say that the pure and sincere faith is signified by gold. Whereof Saint Peter said that the faith of our hearts should be purified. For although there be in us spots and infirmities, yet is faith, by reason of the subject upon which it rests, most pure.

The word of promise, and even Christ himself, is the object of faith, which is purity itself. Therefore, the Lord advises the congregation of Laodicea to be tested by fire, and advises them to hear God’s word and believe it in deed. For the Lord uses the word of being for receiving, hearing, and obeying.

\index{Antichrist}\index{John, Saint}\index{Papacy}For no man should suppose that there is bargaining before God, as there is with men: as though the spiritual gifts of God might be bought for money. This is contrary to the entire Scripture, and especially against the judgment of Saint Peter pronounced against Simon Magus. But this our exposition the Prophet Isaiah approves in the fifty-fifth chapter, where among other things it says, “Come without money, and without price, or exchange.” And immediately following: “Listen carefully, and give your attention.” Therefore, the Roman Antichrist has no claim to this, I mean the Pope, that great merchant, who sells all things in the church, even those things which he does not possess, the greatest deceiver in the world. Moreover, just as it is plainly expressed in Isaiah concerning from whom such graces or gifts are to be bought, so also Christ here says expressly, “I counsel you to buy of me.” Behold, he says of me, not of the Pope, of monks, friars, or priests. For Christ alone has the things which we may require. He alone satisfies, he alone grants those gifts. And therefore he says in the Gospel of Saint John: “Let him who is hungry or thirsty come to me.” To me, I say, let him come. John 4:6 and 7. And Saint Peter says, “Lord, to whom shall we go? You have the words of eternal life.” As though he should say: If we will live, we can go to no other, but unto you. You are the life and fountain of all goodness.

Moreover, the use and benefit of this pure gold—that is, the word of God’s truth and pure faith, tried and most purified—is threefold. First, that you may be rich; secondly, that you may purchase clothing; thirdly, that you may buy the eye salve, to heal the blindness of your eyes. For the word of God and faith in him are the foundation of true piety. Without the word and faith, nothing is sound.

The first fruit is wealth or riches, namely spiritual. For the word and faith are not a false imagination and a vain dream of excellent things. For he who believes the word feels joy in his heart and enjoys spiritual gifts; and possessing Christ through faith possesses all goodness. Whereupon also the Apostle in the first chapter of the first epistle to the Corinthians said: “I give thanks to my God always for you, for the grace of God that is given to you in Christ Jesus, because you are in all things enriched by him, in every word, and in all knowledge, as the testimony of Christ is confirmed in you. In so much that you are not destitute in any gift.” Consider these things well, those who think worldly goods to be true riches. These people shall be judged by the wisdom of God, as it is manifest in the twelfth chapter of Saint Luke. And besides this, those who are destitute of the light of God’s word and lack faith cannot rightly use these earthly riches. Therefore, heavenly riches are the true riches.

The second fruit is the clothing and beautiful apparel with which we are covered, so that our shameful nakedness should not be revealed. Before their fall, our parents were naked, but without any shame or disgrace. After the fall, they were ashamed. Because sin brings shame, and a lack of all good works; and an ungodly lifestyle is a most shameful nakedness.

The Laodiceans were afflicted with this. But Christ, who is learned through the word of truth and perceived by true faith, is the white apparel of the faithful, their righteousness and innocence. He covers all our blemishes, he abolishes our shameful nakedness, decks us with all kinds of virtues, so that we may appear honorable and comely before God in holy conduct. For Christ is the wedding garment. The Apostle counsels us to put on Christ, and that we be clothed with righteousness, temperance, and all goodness. The passages are in Romans 13, Ephesians 4, and Colossians 3. Away with the cowl of our Lady, under which gather for the most part wicked and impenitent persons. The most pure virgin does not cover such [?], she loves righteousness and repentance.

Finally, with this gold is bought an eye salve, which is a medicine for the eyes, as physicians are accustomed to apply to sore and blurred eyes, to prevent blindness. The commandment of the Lord, says David, is bright, giving light to the eyes. Faith also rightly informs human judgment, so that we may rightly discern virtues and vices. The lack of God’s word and true faith brings about blindness.

For all these things the Lord exhorts the people of Laodicea to seek God’s word and believe it faithfully. Thus it will come to pass that, being enriched with all spiritual gifts, they may live a blameless life within the church, may possess Christ, and rightly judge all matters of salvation. And in these things also consists true repentance: forgiveness of sins and amendment of life. \&c.

\index{Hebrews, Epistle to the}But lest they should say, “We hear these things in vain,” as those who have heard before that we shall be spewed out of the Lord’s mouth; yea, and are so sharply rebuked with bitter words and sentences that we are compelled to despair—he anticipates that objection and says, “Whomsoever I love, I rebuke and chasten.” The first word signifies to accuse and reprove openly, which is done with sharper words; the latter is referred to discipline, whereby children are kept in awe with the rod, lest they forget themselves through wantonness. The Lord therefore, alluding to the words of Solomon in chapter 3, signifies that a sharp rebuke or severe chastening is not always a sign that God is angry, but often a token that he is pleased and loves us. Therefore he says, “First, I rebuked you sharply out of love, and so sought your salvation.” Therefore, it is now also a wholesome sign if the preachers rebuke the church with sharp words; and again, it is an unlucky sign if a fox’s tail is stroked over intolerable faults. It is a token of love also, if a man suffers sundry mishaps. Which thing the Apostle discourses at large in chapter 12 to the Hebrews.

Upon these things he infers the substance, and saith: Where you see God so earnestly seeking your salvation, I pray you do not always remain in a murmuring, neither hot nor cold. Be zealous, take unto yourself a fervent zeal to follow and grasp your salvation. For now he sets the fervor of faith, conceived by the word and spirit of God, against this neutrality or lukewarmness. After he adds, and repent, in forsaking your evil conversation, and being of Christ tried gold: that is, purified and purged, whereby you may be rich, be arrayed in white, and may have a medicine wherewith to anoint your eyes, that you may see. To God be glory.
